% =====================================================================
%  Seance 5 -- E. TP B1-B5, quiz de synthese, kbox, preview seance 6
% =====================================================================
\section{TP « Détecter l'intrusion de NOSYS » et synthèse}

\begin{frame}{Le TP : objectif final, une note d'incident}
\begin{columns}[t]
\begin{column}{0.52\textwidth}
\begin{kbox}[Ce que vous produisez à 18 h]
\begin{center}
\file{sessions\_suspectes.csv} + \textbf{note d'incident (1 page)}\\[1mm]
une ligne $=$ (machine, fenêtre 10 min)\\[1mm]
\begin{tabular}{@{}l@{}}
\toprule
echecs, succes, utilisateurs, entropie\\
echecs\_consecutifs, comptes\_sensibles, nuit\\
score\_anomalie (IsolationForest)\\
pred\_attaque (forêt supervisée)\\
\bottomrule
\end{tabular}
\end{center}
\end{kbox}
\begin{ebox}[Chiffres attendus (mesurés)]
\textbf{10\,000 sessions} dont \textbf{12 d'attaque} (0,12 \%) ;
top 30 de l'IF $=$ les 12 sessions d'attaque aux \textbf{rangs 1--12},
précision@30 $= 0{,}40$, rappel $= 1{,}00$ ; forêt supervisée :
précision $=$ rappel $= 1{,}00$ en test stratifié.
\end{ebox}
\end{column}
\begin{column}{0.44\textwidth}
\begin{wbox}[Barème (20 points)]
\begin{itemize}
\item B1 parsing + agrégation : \textbf{4 pts}
\item B2 features + entropie : \textbf{4 pts}
\item B3 IsolationForest + précision@30 : \textbf{5 pts}
\item B4 forêt supervisée + métriques : \textbf{4 pts}
\item B5 livrable + note d'incident : \textbf{3 pts}
\end{itemize}
Graine \textbf{5534} partout, \textbf{2 threads} max. Le code doit
s'exécuter de bout en bout sur \file{nosys\_logs.csv}.
\end{wbox}
\end{column}
\end{columns}
\end{frame}

\begin{frame}[fragile]{TP B1 --- parser et agréger (30 min)}
\begin{codebox}[B1. De 100 000 lignes à 10 000 sessions]
\begin{lstlisting}
import pandas as pd

L = pd.read_csv("data/nosys_logs.csv", parse_dates=["horodatage"])
L["echec"]  = (~L["auth_reussie"].astype(bool)).astype(int)
L["succes"] = 1 - L["echec"]
L["fenetre"] = L["horodatage"].dt.floor("10min")

g = L.groupby(["machine", "fenetre"])
S = pd.DataFrame({
    "n_evenements": g.size(),
    "echecs":  g["echec"].sum(),
    "succes":  g["succes"].sum(),
    "utilisateurs": g["utilisateur"].nunique(),
})
S["taux_echec"] = S["echecs"] / S["n_evenements"]
print(len(S))            # 10000 sessions
print(S["echecs"].max()) # 15  (la rafale de l'attaque)
\end{lstlisting}
\end{codebox}
\vspace{1mm}
\begin{solbox}[Vérification attendue]
10\,000 sessions ; max \py{echecs} $= 15$ (bdd-principal, 03 h 10) ;
moyenne d'échecs par session $\approx 0{,}48$ --- l'attaque dépasse la
normale d'un facteur 20 sur sa feature principale.
\end{solbox}
\end{frame}

\begin{frame}[fragile]{TP B2 --- features d'attaque et entropie (30 min)}
\begin{codebox}[B2. Entropie de Shannon des comptes testés]
\begin{lstlisting}
from scipy.stats import entropy

def entropie_comptes(u):          # u = liste des utilisateurs
    comptage = pd.Series(list(u)).value_counts().values
    return entropy(comptage, base=2)

S["entropie"] = g["utilisateur"].apply(entropie_comptes)

# echecs consecutifs max (la rafale), ordre horaire
def max_run(d):
    e = d.sort_values("horodatage")["echec"].values
    best = cur = 0
    for v in e:
        cur = cur + 1 if v else 0
        best = max(best, cur)
    return best

S["echecs_consecutifs"] = g.apply(max_run, include_groups=False)
SENS = {"root", "admin", "inconnu-1", "inconnu-2"}
S["comptes_sensibles"] = g.apply(
    lambda d: int(d.loc[d["echec"]==1,
                        "utilisateur"].isin(SENS).sum()),
    include_groups=False)
S["nuit"] = (S.index.get_level_values("fenetre").hour < 6)
\end{lstlisting}
\end{codebox}
\vspace{1mm}
\begin{center}\small
Vérification : entropie moyenne attaque $\approx 1{,}73$ bits vs
$2{,}22$ bits sur le reste ; \textbf{100 \%} des échecs d'attaque portent
sur des comptes sensibles.
\end{center}
\end{frame}

\begin{frame}[fragile]{TP B3 --- IsolationForest et précision@30 (30 min)}
\footnotesize%
\begin{codebox}[B3. Classer, puis compter le recouvrement]
\begin{lstlisting}
from sklearn.ensemble import IsolationForest

FEATS = ["echecs", "taux_echec", "utilisateurs", "entropie",
         "echecs_consecutifs", "comptes_sensibles", "nuit"]
X = S[FEATS].values

iso = IsolationForest(n_estimators=200, contamination=0.005,
                      random_state=5534, n_jobs=2)
S["score"] = -iso.fit(X).decision_function(X)
top30 = S.sort_values("score", ascending=False).head(30)

# verite terrain : fenetre d'attaque reelle
att = (top30["machine"] == "bdd-principal") & \
      top30["fenetre"].between("2026-10-05 02:00",
                               "2026-10-05 04:00")
print("precision@30 :", att.mean())   # 0.40
print("rappel :", ...)                 # 12/12 = 1.0
\end{lstlisting}
\end{codebox}
\vspace{1mm}
\begin{center}\small
Attendu : les 12 sessions d'attaque occupent les \textbf{rangs 1 à 12}.
Discutez : pourquoi précision@10 $= 1{,}0$ mais précision@30 $= 0{,}4$ ?
\end{center}
\end{frame}

\begin{frame}[fragile]{TP B4 --- forêt supervisée ; B5 --- livrable (30 min)}
\begin{columns}[t]
\begin{column}{0.52\textwidth}
\begin{codebox}[B4. Supervisé sur l'étiquette session-attaque]
\begin{lstlisting}
from sklearn.ensemble import RandomForestClassifier
from sklearn.model_selection import \
    train_test_split

S["attaque"] = att_reelle(S)   # fenetre 2-4h bdd
Xtr, Xte, ytr, yte = train_test_split(
    S[FEATS], S["attaque"], test_size=0.5,
    stratify=S["attaque"], random_state=5534)
rf = RandomForestClassifier(
    n_estimators=300, class_weight="balanced",
    random_state=5534, n_jobs=2).fit(Xtr, ytr)
print(precision_score(yte, rf.predict(Xte)))  # 1.0
print(recall_score(yte, rf.predict(Xte)))     # 1.0
\end{lstlisting}
\end{codebox}
\end{column}
\begin{column}{0.44\textwidth}
\begin{ebox}[B5. Livrable (10 min)]
\begin{itemize}
\item \file{sessions\_suspectes.csv} : le top 30 avec score et prédiction.
\item \textbf{Note d'incident (1 page)} : ce qui s'est passé (brute force,
bdd-principal, 05/10 02 h -- 04 h), \textbf{preuve chiffrée} (125 échecs,
12 sessions, comptes ciblés, 0 succès), \textbf{action recommandée}
(ex. blocage IP + verrouillage progressif + rotation des mots de passe).
\end{itemize}
\end{ebox}
\begin{wbox}[Attention]
\py{att\_reelle} ne doit servir \textbf{qu'à l'évaluation} --- l'utiliser
comme feature serait une fuite de données (règle anti-fuite de la séance 2).
\end{wbox}
\end{column}
\end{columns}
\end{frame}

\begin{frame}{Quiz de synthèse --- 5 QCM (10 min)}
\begin{qbox}[Séance 5, tout le programme]
\begin{enumerate}
\item Le premier réflexe face à 100\,000 lignes de logs est de :\\
(a) tout lire \quad (b) agréger en sessions et construire des features \quad
(c) acheter un SIEM
\item La granularité « session » utilisée aujourd'hui est :\\
(a) (utilisateur, jour) \quad (b) (machine, fenêtre 10 min) \quad (c) (message, niveau)
\item IsolationForest a placé les 12 sessions d'attaque :\\
(a) aux rangs 1--12 \quad (b) aux rangs 100--112 \quad (c) nulle part
\item Le superviseur bat le non-supervisé surtout parce qu'il :\\
(a) va plus vite \quad (b) connaît la frontière attaque/normal apprise des
incidents passés \quad (c) n'a pas d'hyperparamètres
\item La meilleure défense contre la fatigue d'alerte est :\\
(a) plus d'alertes \quad (b) classer et regrouper pour ne montrer que le
haut du classement \quad (c) supprimer les règles
\end{enumerate}
\end{qbox}
\vspace{1mm}
\begin{solbox}[Corrigé]
\textbf{1 : b} \quad \textbf{2 : b} \quad \textbf{3 : a} (mesuré, graine 5534) \quad
\textbf{4 : b} \quad \textbf{5 : b}
\end{solbox}
\end{frame}

\begin{frame}{Ce qu'il faut retenir --- et la suite}
\begin{columns}[t]
\begin{column}{0.52\textwidth}
\begin{kbox}[Les 5 idées de la séance]
\begin{enumerate}
\item La signature d'une attaque vit dans l'\textbf{agrégation}, pas dans
la ligne : sessions (machine $\times$ 10 min) et features d'attaque.
\item \textbf{Combiner des signaux faibles} : échecs + entropie + comptes
sensibles + nuit.
\item IsolationForest \textbf{sans étiquette} a retrouvé les 12 sessions
d'attaque aux rangs 1--12 (précision@30 $= 0{,}40$, hasard : 0,0012).
\item Le supervisé (forêt, étiquettes d'incidents) : précision $=$ rappel
$= 1{,}00$ en test --- mais méfiance devant le 100 \%.
\item Limites à garder en tête : adversarial, fatigue d'alerte,
explicabilité (SHAP), vie privée ; fiabilité : dérive, données, redondance.
\end{enumerate}
\end{kbox}
\end{column}
\begin{column}{0.44\textwidth}
\begin{ebox}[Prochaine séance --- Séance 6]
\textbf{MLOps : faire vivre les modèles}
\begin{itemize}
\item Exposer un détecteur en \textbf{API} (endpoint de scoring).
\item \textbf{Réentraînement} automatique : déclencheurs, fenêtres,
validation avant mise en production.
\item Surveillance de bout en bout : données, dérive, scores --- le
pipeline complet de NOSYS, de la collecte à l'alerte.
\end{itemize}
À préparer : relire la séance 3 (forêt, validation croisée) --- la séance 6
met le \emph{modèle} en production, pas seulement le code.
\end{ebox}
\end{column}
\end{columns}
\end{frame}
